% GENERATED FILE. Edit canonical lessons, then run npm run generate:books.
% canonical-source-hash: fnv1a64:d45891542c18a339
% canonical-lessons: SA-C531-vanchati, SA-C531-manorathah, SA-C531-aradhayati, SA-C531-prarthana, SA-C531-kamana

\chapter{Wishes and Prayers}
\label{ch:sa-wishes-and-prayers}
\hlchaptermodality{531}

\begin{chapteropening}
\textbf{By the end of this chapter:} \emph{I can say what I wish for and what will come of it.}

Complete the last lesson of chapter 531: I can say what I wish for and what will come of it.
\end{chapteropening}

\section[\texorpdfstring{\sk{वाञ्छति} (vāñchati)}{वाञ्छति (vāñchati)}]{\sk{वाञ्छति} (vāñchati) — yearns for}

\label{lesson:SA-C531-vanchati}



\begin{quote}
Before the new one: say the Sanskrit for a sherbet, then the Sanskrit for stale.
\end{quote}

\subsection*{You'll want to know: \sk{वाञ्छति}}
\textbf{\sk{वाञ्छति}} — \emph{vāñchati} — \textquotedblleft{}yearns for\textquotedblright{}.

\begin{quote}
\sk{यदि} \sk{सः} \sk{सुखं} \sk{वाञ्छति}, \sk{तर्हि} \sk{सः} \sk{परिश्रमं} \sk{करोतु।}
\end{quote}

\begin{quote}
\emph{yadi saḥ sukhaṃ vāñchati, tarhi saḥ pariśramaṃ karotu.}
\end{quote}

\begin{quote}
If he wants happiness, let him work hard.
\end{quote}

Past (he …): \sk{सः} \sk{सुखं} \sk{वाञ्छितवान्}, \emph{saḥ sukhaṃ vāñchitavān}.

\begin{cousinweb}
\textbf{\sk{वाञ्छति}} (\emph{vāñchati}) comes from the root √vāñch, \textquotedblleft{}to desire, to wish\textquotedblright{}; MW compares German \emph{wünschen} and English \emph{wish} (MW).
\end{cousinweb}

\begin{grammarlens}[title={What you've built}]
One more word for saying what happens if.
\end{grammarlens}

\subsection*{Your turn}
\begin{itemize}
\raggedright
  \item \emph{Say it:} \emph{vāñchati}
  \item \emph{Say it:} \emph{vāñchati}, once more
  \item \emph{Say it:} \emph{yadi saḥ sukhaṃ vāñchati, tarhi saḥ pariśramaṃ karotu.} — If he wants happiness, let him work hard.
  \item \emph{Recall:} say the Sanskrit for a sherbet, then the Sanskrit for stale, then say \emph{vāñchati} again
\end{itemize}

\begin{tcolorbox}[breakable,colback=teal!4,colframe=teal!35!black,title={Before you move on}]
What does \sk{वाञ्छति} mean? (\textquotedblleft{}Yearns for\textquotedblright{}.) Say it once more.
\end{tcolorbox}
\section[\texorpdfstring{\sk{मनोरथः} (manorathaḥ)}{मनोरथः (manorathaḥ)}]{\sk{मनोरथः} (manorathaḥ) — a heart's wish}

\label{lesson:SA-C531-manorathah}



\begin{quote}
Before the new one: say the Sanskrit for stale, then the Sanskrit for yearns for.
\end{quote}

\subsection*{You'll want to know: \sk{मनोरथः}}
\textbf{\sk{मनोरथः}} — \emph{manorathaḥ} — \textquotedblleft{}a heart's wish\textquotedblright{}.

\begin{quote}
\sk{यदि} \sk{मनोरथः} \sk{दृढः} \sk{अस्ति}, \sk{तर्हि} \sk{मार्गः} \sk{अपि} \sk{भवति।}
\end{quote}

\begin{quote}
\emph{yadi manorathaḥ dṛḍhaḥ asti, tarhi mārgaḥ api bhavati.}
\end{quote}

\begin{quote}
If your wish is strong, a way appears too.
\end{quote}

A masculine noun.

\begin{cousinweb}
\textbf{\sk{मनोरथः}} (\emph{manorathaḥ}) is built from \emph{manas}, \textquotedblleft{}the mind, the heart\textquotedblright{}, and \emph{ratha}, which MW here takes as \textquotedblleft{}joy\textquotedblright{}: a heart's joy, a wish, a desire (MW).
\end{cousinweb}

\begin{grammarlens}[title={What you've built}]
One more word for saying what happens if.
\end{grammarlens}

\subsection*{Your turn}
\begin{itemize}
\raggedright
  \item \emph{Say it:} \emph{manorathaḥ}
  \item \emph{Say it:} \emph{manorathaḥ}, once more
  \item \emph{Say it:} \emph{yadi manorathaḥ dṛḍhaḥ asti, tarhi mārgaḥ api bhavati.} — If your wish is strong, a way appears too.
  \item \emph{Recall:} say the Sanskrit for stale, then the Sanskrit for yearns for, then say \emph{manorathaḥ} again
\end{itemize}

\begin{tcolorbox}[breakable,colback=teal!4,colframe=teal!35!black,title={Before you move on}]
What does \sk{मनोरथः} mean? (\textquotedblleft{}A heart's wish\textquotedblright{}.) Say it once more.
\end{tcolorbox}
\section[\texorpdfstring{\sk{आराधयति} (ārādhayati)}{आराधयति (ārādhayati)}]{\sk{आराधयति} (ārādhayati) — propitiates; honours}

\label{lesson:SA-C531-aradhayati}



\begin{quote}
Before the new one: say the Sanskrit for yearns for, then the Sanskrit for a heart's wish.
\end{quote}

\subsection*{You'll want to know: \sk{आराधयति}}
\textbf{\sk{आराधयति}} — \emph{ārādhayati} — \textquotedblleft{}propitiates; honours\textquotedblright{}.

\begin{quote}
\sk{यदि} \sk{त्वं} \sk{गुरुम्} \sk{आराधयसि}, \sk{तर्हि} \sk{सः} \sk{प्रसन्नः} \sk{भविष्यति।}
\end{quote}

\begin{quote}
\emph{yadi tvaṃ gurum ārādhayasi, tarhi saḥ prasannaḥ bhaviṣyati.}
\end{quote}

\begin{quote}
If you honour your teacher, he will be pleased.
\end{quote}

Past (he …): \sk{सः} \sk{गुरुम्} \sk{आराधितवान्}, \emph{saḥ gurum ārādhitavān}.

\begin{cousinweb}
\textbf{\sk{आराधयति}} (\emph{ārādhayati}) is built on \emph{ā} and the causative of the root √rādh, \textquotedblleft{}to succeed\textquotedblright{}; MW gives \emph{ārādhayati} as to conciliate, to propitiate, to honour, to worship (MW).
\end{cousinweb}

\begin{grammarlens}[title={What you've built}]
One more word for saying what happens if.
\end{grammarlens}

\subsection*{Your turn}
\begin{itemize}
\raggedright
  \item \emph{Say it:} \emph{ārādhayati}
  \item \emph{Say it:} \emph{ārādhayati}, once more
  \item \emph{Say it:} \emph{yadi tvaṃ gurum ārādhayasi, tarhi saḥ prasannaḥ bhaviṣyati.} — If you honour your teacher, he will be pleased.
  \item \emph{Recall:} say the Sanskrit for yearns for, then the Sanskrit for a heart's wish, then say \emph{ārādhayati} again
\end{itemize}

\begin{tcolorbox}[breakable,colback=teal!4,colframe=teal!35!black,title={Before you move on}]
What does \sk{आराधयति} mean? (\textquotedblleft{}Propitiates; honours\textquotedblright{}.) Say it once more.
\end{tcolorbox}
\section[\texorpdfstring{\sk{प्रार्थना} (prārthanā)}{प्रार्थना (prārthanā)}]{\sk{प्रार्थना} (prārthanā) — a prayer; a request}

\label{lesson:SA-C531-prarthana}



\begin{quote}
Before the new one: say the Sanskrit for a heart's wish, then the Sanskrit for propitiates.
\end{quote}

\subsection*{You'll want to know: \sk{प्रार्थना}}
\textbf{\sk{प्रार्थना}} — \emph{prārthanā} — \textquotedblleft{}a prayer; a request\textquotedblright{}.

\begin{quote}
\sk{यदि} \sk{प्रार्थना} \sk{शुद्धा} \sk{अस्ति}, \sk{तर्हि} \sk{मनः} \sk{शान्तं} \sk{भवति।}
\end{quote}

\begin{quote}
\emph{yadi prārthanā śuddhā asti, tarhi manaḥ śāntaṃ bhavati.}
\end{quote}

\begin{quote}
If the prayer is pure, the mind grows calm.
\end{quote}

A feminine noun.

\begin{cousinweb}
\textbf{\sk{प्रार्थना}} (\emph{prārthanā}) is built on \emph{pra} and the root √arth, \textquotedblleft{}to ask for\textquotedblright{}; MW gives \emph{prārthanā} as a request, a petition, and a prayer in worship (MW).
\end{cousinweb}

\begin{grammarlens}[title={What you've built}]
One more word for saying what happens if.
\end{grammarlens}

\subsection*{Your turn}
\begin{itemize}
\raggedright
  \item \emph{Say it:} \emph{prārthanā}
  \item \emph{Say it:} \emph{prārthanā}, once more
  \item \emph{Say it:} \emph{yadi prārthanā śuddhā asti, tarhi manaḥ śāntaṃ bhavati.} — If the prayer is pure, the mind grows calm.
  \item \emph{Recall:} say the Sanskrit for a heart's wish, then the Sanskrit for propitiates, then say \emph{prārthanā} again
\end{itemize}

\begin{tcolorbox}[breakable,colback=teal!4,colframe=teal!35!black,title={Before you move on}]
What does \sk{प्रार्थना} mean? (\textquotedblleft{}A prayer; a request\textquotedblright{}.) Say it once more.
\end{tcolorbox}
\section[\texorpdfstring{\sk{कामना} (kāmanā)}{कामना (kāmanā)}]{\sk{कामना} (kāmanā) — a desire}

\label{lesson:SA-C531-kamana}



\begin{quote}
Before the new one: say the Sanskrit for propitiates, then the Sanskrit for a prayer.
\end{quote}

\subsection*{You'll want to know: \sk{कामना}}
\textbf{\sk{कामना}} — \emph{kāmanā} — \textquotedblleft{}a desire\textquotedblright{}.

\begin{quote}
\sk{यदि} \sk{तव} \sk{कामना} \sk{उचिता} \sk{अस्ति}, \sk{तर्हि} \sk{सा} \sk{पूर्णा} \sk{भविष्यति।}
\end{quote}

\begin{quote}
\emph{yadi tava kāmanā ucitā asti, tarhi sā pūrṇā bhaviṣyati.}
\end{quote}

\begin{quote}
If your desire is right, it will be fulfilled.
\end{quote}

A feminine noun.

\begin{cousinweb}
\textbf{\sk{कामना}} (\emph{kāmanā}) comes from the root √kam, \textquotedblleft{}to wish, to love\textquotedblright{}; MW gives \emph{kāmanā} as a wish, a desire (MW).
\end{cousinweb}

\begin{grammarlens}[title={What you've built}]
One more word for saying what happens if.
\end{grammarlens}

\subsection*{Your turn}
\begin{itemize}
\raggedright
  \item \emph{Say it:} \emph{kāmanā}
  \item \emph{Say it:} \emph{kāmanā}, once more
  \item \emph{Say it:} \emph{yadi tava kāmanā ucitā asti, tarhi sā pūrṇā bhaviṣyati.} — If your desire is right, it will be fulfilled.
  \item \emph{Recall:} say the Sanskrit for propitiates, then the Sanskrit for a prayer, then say \emph{kāmanā} again
\end{itemize}

\begin{tcolorbox}[breakable,colback=teal!4,colframe=teal!35!black,title={Before you move on}]
What does \sk{कामना} mean? (\textquotedblleft{}A desire\textquotedblright{}.) Say it once more.
\end{tcolorbox}
